% ===== Part I, section 1: setting and harmonic preliminaries; section 2: Target 3 =====
\section{Harmonic preliminaries}\label{sec:prelim}

Throughout, $X,X_1,\dots,X_n$ are independent and uniform on $\Sph$, $d\ge3$, and
\[
K(x,y)=\exp\Big(-\frac{\|x-y\|^2}{2h^2}\Big)=e^{-\gamma}e^{\gamma\langle x,y\rangle},\qquad \gamma=h^{-2},
\]
with $h>0$ fixed. We write $\mathcal H_{d,s}$ for the space of spherical harmonics of degree $s$ on
$\Sph$, $N_s=N_{d,s}=\dim\mathcal H_{d,s}$, and $\Pi_s(x,y)=\sum_{a=1}^{N_s}\psi_{s,a}(x)\psi_{s,a}(y)$
for its reproducing kernel with respect to the uniform probability measure, so that
$\Pi_s(x,y)=N_sP_s(\langle x,y\rangle)$ with $P_s$ the zonal harmonic normalised by $P_s(1)=1$
\cite{AtkinsonHan}. The addition theorem gives $\Pi_s(x,x)=N_s$, and orthogonality of distinct degrees
gives, for every $x,z\in\Sph$,
\begin{equation}
\E_Y\big[\Pi_s(x,Y)\Pi_t(Y,z)\big]=\delta_{st}\,\Pi_s(x,z),\qquad \E_Y\Pi_s(x,Y)=0\quad(s\ge1).
\label{eq:repro}
\end{equation}
The Mercer expansion of the kernel is $K(x,y)=\sum_{s\ge0}\lambda_s\Pi_s(x,y)$ with $\lambda_s=\lambda_{d,s}>0$
\cite{MinhNiyogiYao}; since $K(x,x)=1$, the trace masses $\tau_s=N_s\lambda_s$ satisfy $\sum_s\tau_s=1$.
For a sample we write $K=(K(X_i,X_j))_{ij}$ for the Gram matrix, $\Pi_s(X,X)=(\Pi_s(X_i,X_j))_{ij}$, and
$K_s=\lambda_s\Pi_s(X,X)$, so that $K=\sum_sK_s$. Each $\Pi_s(X,X)=\Phi_s\Phi_s^\top$ with
$\Phi_s=(\psi_{s,a}(X_i))_{i,a}\in\R^{n\times N_s}$, and $\|\Phi_s\|_F^2=\sum_i\Pi_s(X_i,X_i)=nN_s$.

\begin{lemma}[harmonic multiplicities]\label{lem:mult}
For $d\ge3$ and $s\ge0$,
\[
N_{d,s}=\binom{d+s-1}{s}-\binom{d+s-3}{s-2}=\frac{(2s+d-2)\,(s+d-3)!}{s!\,(d-2)!},
\]
\[
\rho_s:=\frac{N_{d,s+1}}{N_{d,s}}=\frac{(d+2s)(d+s-2)}{(s+1)(d+2s-2)} .
\]
The ratio $\rho_s$ is at least one and is non-increasing in $s$. For $s\le d$ one has
$\tfrac34\binom{d+s-1}{s}\le N_{d,s}\le\binom{d+s-1}{s}$, and for $s^2=o(d)$,
$N_{d,s}=\frac{d^s}{s!}\,e^{O(s^2/d)}$.
\end{lemma}

\begin{proof}
The closed forms are classical \cite{AtkinsonHan}. From the product form,
$\rho_s=\frac{d+s-2}{s+1}\cdot\frac{d+2s}{d+2s-2}$; both factors are non-increasing in $s$, and
$\frac{d+s-2}{s+1}\ge1$ for $d\ge3$. The ratio of the two binomial coefficients is
$\frac{s(s-1)}{(d+s-1)(d+s-2)}\le\big(\frac{s}{d+s}\big)^2\le\frac14$ for $s\le d$, which gives the two-sided
bound. For $s^2=o(d)$, $\binom{d+s-1}{s}=\frac{d^s}{s!}\prod_{j=1}^{s-1}(1+j/d)=\frac{d^s}{s!}e^{O(s^2/d)}$
and the subtracted term is smaller by the factor $O(s^2/d^2)$.
\end{proof}

\begin{lemma}[Mercer eigenvalues of the Gaussian kernel]\label{lem:mercer}
With $(d)_{m}:=d(d+2)\cdots(d+2m-2)$ and $(d)_0=1$,
\begin{equation}
\lambda_{d,s}=e^{-\gamma}\sum_{r\ge0}\frac{\gamma^{s+2r}}{2^r\,r!\,(d)_{s+r}}
=e^{-\gamma}\,\Gamma\!\big(\tfrac d2\big)\Big(\frac2\gamma\Big)^{d/2-1}I_{s+d/2-1}(\gamma),
\label{eq:mercer}
\end{equation}
where $I_\nu$ is the modified Bessel function of the first kind. Consequently, for all $s\ge0$,
\begin{equation}
\frac{\lambda_{d,s+1}}{\lambda_{d,s}}\le\frac{\gamma}{d+2s},\qquad
\lambda_{d,s}\ge\frac{e^{-\gamma}\gamma^s}{(d)_s},\qquad
\tau_s\le e^{-\gamma}\sum_{m\ge s}\frac{\gamma^m}{m!}.
\label{eq:mercerbounds}
\end{equation}
In particular the eigenvalues are strictly decreasing in $s$ as soon as $d>\gamma$.
\end{lemma}

\begin{proof}
Write $K(x,y)=e^{-\gamma}\sum_m\frac{\gamma^m}{m!}\langle x,y\rangle^m$. The Funk--Hecke formula
\cite{AtkinsonHan} gives the degree-$s$ coefficient of $t\mapsto t^m$ on $\Sph$; equivalently, for
$G\sim N(0,I_d)$ and the homogeneous harmonic extension $H_s(ru)=r^sY_s(u)$ of $Y_s\in\mathcal H_{d,s}$,
the identity $\E[H_s(G)e^{t\langle x,G\rangle-t^2/2}]=t^sH_s(x)$ and the comparison of the coefficient of
$t^{s+2r}$ give $\E[\langle x,G\rangle^{s+2r}H_s(G)]=\frac{(s+2r)!}{2^rr!}H_s(x)$, while
$\E\|G\|^{2(s+r)}=(d)_{s+r}$. Dividing out the radial moments yields the series in \eqref{eq:mercer};
the Bessel form is the series $I_\nu(\gamma)=\sum_r\frac{(\gamma/2)^{\nu+2r}}{r!\,\Gamma(\nu+r+1)}$
\cite[10.25.2]{DLMF} rewritten with $\nu=s+d/2-1$. The first two bounds in \eqref{eq:mercerbounds} are
termwise: the $r$-th term of $\lambda_{s+1}$ equals $\frac{\gamma}{d+2(s+r)}$ times the $r$-th term of
$\lambda_s$, and $\frac{\gamma}{d+2(s+r)}\le\frac{\gamma}{d+2s}$; the second is the $r=0$ term. For the
third, expand $\langle x,y\rangle^m=\sum_{s\le m}c_{m,s}P_s(\langle x,y\rangle)$ with $c_{m,s}\ge0$
(the Gegenbauer coefficients of a monomial are non-negative) and $\sum_sc_{m,s}=1$ by evaluating at
$\langle x,y\rangle=1$; hence $\tau_s=e^{-\gamma}\sum_{m\ge s}\frac{\gamma^m}{m!}c_{m,s}\le
e^{-\gamma}\sum_{m\ge s}\frac{\gamma^m}{m!}$.
\end{proof}

\begin{lemma}[hypercontractivity at growing degree]\label{lem:hyper}
For every $Y_s\in\mathcal H_{d,s}$, every $s\ge0$ and every $p\ge2$,
$\|Y_s\|_{L^p(\Sph)}\le(p-1)^{s/2}\|Y_s\|_{L^2(\Sph)}$. In particular, with $r_s:=\Pi_s/N_s$ and $z\in\Sph$
fixed, $\|r_s(X,z)\|_{L^{2q}}\le(2q-1)^{s/2}N_s^{-1/2}$ for every $q\ge1$.
\end{lemma}

\begin{proof}
Let $G\sim N(0,I_d)$ and $H_s(G)=\|G\|^sY_s(G/\|G\|)$, a homogeneous harmonic polynomial of degree $s$
and hence an element of the $s$-th Gaussian chaos. Nelson's hypercontractive inequality
\cite{Nelson,Janson} gives $\|H_s(G)\|_p\le(p-1)^{s/2}\|H_s(G)\|_2$. The radius and the direction of $G$
are independent, so $\|H_s(G)\|_p=\|\,\|G\|^s\|_p\,\|Y_s\|_{L^p}$ and $\|H_s(G)\|_2=\|\,\|G\|^s\|_2\,\|Y_s\|_{L^2}$,
and $\|\,\|G\|^s\|_p\ge\|\,\|G\|^s\|_2$ for $p\ge2$. Dividing gives the first claim; the second follows
from $\|r_s(\cdot,z)\|_{L^2}^2=\Pi_s(z,z)/N_s^2=1/N_s$ by \eqref{eq:repro}.
\end{proof}

\section{Target 3: the exponential-sample lower bound}\label{sec:proofT3}

Here $n=2^d$, the signal is zero, $Y_i=\varepsilon_i$ with $\varepsilon\sim N(0,\sigma^2I_n)$, and
$\hat f(x)=k(x)^\top K^{-1}\varepsilon$ is the minimum-RKHS-norm interpolant, $k(x)=(K(x,X_i))_i$. The Gaussian
kernel is strictly positive definite, so $K$ is invertible almost surely. The excess risk is
$\mathcal E_{n,d}=\E_{X_0}[\hat f(X_0)^2]$, and we write $V(X)=\E_\varepsilon[\mathcal E_{n,d}\mid X]$.

\begin{theorem}[Target 3]\label{thm:T3}
For every fixed $h>0$,
\[
\liminf_{d\to\infty}\ \frac{\E\,\mathcal E_{2^d,d}}{\sigma^2}\ \ge\ \frac5{36},
\]
the expectation being over the design and the noise.
\end{theorem}

The proof has three steps: an exact reduction of $V(X)$ to the leverage of a block of harmonic degrees,
a lower bound on the expected leverage of a block, and the location of the last block whose dimension
does not exceed $n$.

\begin{lemma}[reduction to a block]\label{lem:reduction}
Let $G=\E_{X_0}[k(X_0)k(X_0)^\top]=\sum_s\lambda_sK_s$, $t_s=\Tr(K^{-1}K_s)$ and $q_s=\Tr(K^{-1}K_sK^{-1})$.
Then $V(X)/\sigma^2=\Tr(K^{-1}GK^{-1})=\sum_s\lambda_sq_s$, $\sum_st_s=n$, $t_s^2\le nN_s\lambda_sq_s$ for
every $s$, and for every finite set $B$ of degrees, with $T_B=\sum_{s\in B}t_s$ and $D_B=\sum_{s\in B}N_s$,
\begin{equation}
\frac{V(X)}{\sigma^2}\ \ge\ \frac{T_B^2}{nD_B}.
\label{eq:blockreduction}
\end{equation}
\end{lemma}

\begin{proof}
$\E_\varepsilon\hat f(x)^2=\sigma^2k(x)^\top K^{-2}k(x)$, so $V(X)/\sigma^2=\Tr(K^{-1}GK^{-1})$, and by
\eqref{eq:repro},
\[
G_{ij}=\sum_{s,t}\lambda_s\lambda_t\,\E_{X_0}\big[\Pi_s(X_i,X_0)\Pi_t(X_0,X_j)\big]=\sum_s\lambda_s^2\Pi_s(X_i,X_j),
\]
i.e.\ $G=\sum_s\lambda_sK_s$; hence $V/\sigma^2=\sum_s\lambda_sq_s$, and
$\sum_st_s=\Tr(K^{-1}K)=n$. Since $K_s=\lambda_s\Phi_s\Phi_s^\top$, we have
$t_s=\lambda_s\langle K^{-1}\Phi_s,\Phi_s\rangle_F$ and $q_s=\lambda_s\|K^{-1}\Phi_s\|_F^2$, so by
Cauchy--Schwarz
\[
t_s\le\lambda_s\|K^{-1}\Phi_s\|_F\,\|\Phi_s\|_F=\sqrt{\lambda_sq_s}\,\sqrt{nN_s}.
\]
Summing over $B$ and applying Cauchy--Schwarz once more,
$T_B\le\sqrt{nD_B}\big(\sum_{s\in B}\lambda_sq_s\big)^{1/2}\le\sqrt{nD_B\,V(X)/\sigma^2}$.
\end{proof}

\begin{lemma}[expected leverage of a block]\label{lem:leverage}
Let $d>\gamma$ and $L\ge0$, and put $D=\sum_{s\le L}N_s$, $\lambda=\lambda_{d,L}$, $\tau_{>L}=\sum_{s>L}\tau_s$
and $T_{\le L}=\sum_{s\le L}t_s$. Then
\begin{equation}
\E\,T_{\le L}\ \ge\ \frac{nD}{D+n-1+\tau_{>L}/\lambda}.
\label{eq:leverage}
\end{equation}
\end{lemma}

\begin{proof}
Put $A=\sum_{s\le L}K_s$, $C=K-A=\sum_{s>L}K_s\succeq0$, and $A_w=\lambda\sum_{s\le L}\Pi_s(X,X)$. By
Lemma~\ref{lem:mercer} the eigenvalues decrease in $s$, so $\lambda_s\ge\lambda$ for $s\le L$ and
$0\preceq A_w\preceq A$. For $\epsilon>0$ let $K_\epsilon=K+\epsilon I$ and $C_\epsilon=A_w+C+\epsilon I$; then
$K_\epsilon\succeq C_\epsilon\succ0$, so $K_\epsilon^{-1}\preceq C_\epsilon^{-1}$ by operator monotonicity of the
inverse. Since $A+C+\epsilon I=K_\epsilon$,
\[
\Tr(K_\epsilon^{-1}A)=n-\Tr\big(K_\epsilon^{-1}(C+\epsilon I)\big)\ \ge\ n-\Tr\big(C_\epsilon^{-1}(C+\epsilon I)\big)
=\Tr(C_\epsilon^{-1}A_w),
\]
where the inequality uses $C+\epsilon I\succeq0$. Cauchy--Schwarz for the Hilbert--Schmidt inner product,
applied to $C_\epsilon^{-1/2}A_w^{1/2}$ and $C_\epsilon^{1/2}A_w^{1/2}$, gives
$(\Tr A_w)^2\le\Tr(C_\epsilon^{-1}A_w)\Tr(C_\epsilon A_w)$. By the addition theorem $\Tr A_w=n\lambda D$
pathwise. Letting $\epsilon\downarrow0$ (the Gram matrix is invertible almost surely),
\begin{equation}
T_{\le L}=\Tr(K^{-1}A)\ \ge\ \frac{(n\lambda D)^2}{\Tr(A_w^2)+\Tr(CA_w)}.
\label{eq:pathwise}
\end{equation}
The denominator has exact expectations. By \eqref{eq:repro}, for independent $X,Y$,
$\E[\Pi_{\le L}(X,Y)^2]=\sum_{s\le L}\E_Y\Pi_s(X,Y)^2=\sum_{s\le L}\Pi_s(X,X)=D$ where
$\Pi_{\le L}=\sum_{s\le L}\Pi_s$, while $\Pi_{\le L}(X,X)=D$; hence
$\E\Tr(A_w^2)=\lambda^2\big(nD^2+n(n-1)D\big)$. Likewise $\E[C(X,Y)A_w(X,Y)]=0$ for $X\ne Y$ because the
two kernels expand over disjoint sets of degrees, while $C(X,X)=\tau_{>L}$ and $A_w(X,X)=\lambda D$; hence
$\E\Tr(CA_w)=n\tau_{>L}\lambda D$. No independence between $C$ and $A_w$ is used here: both are functions of
the same design, and the off-diagonal expectation vanishes by orthogonality alone. Since $z\mapsto1/z$ is
convex on $(0,\infty)$, Jensen's inequality applied to \eqref{eq:pathwise} gives
\[
\E\,T_{\le L}\ \ge\ \frac{(n\lambda D)^2}{\lambda^2(nD^2+n(n-1)D)+n\tau_{>L}\lambda D}
=\frac{nD}{D+n-1+\tau_{>L}/\lambda}. \qedhere
\]
\end{proof}

\begin{lemma}[the last subcritical block]\label{lem:window}
Let $n=2^d$ and $L_-=L_-(d)=\max\{L:\sum_{s\le L}N_{d,s}\le n\}$, $D_-=\sum_{s\le L_-}N_{d,s}$ and
$y_d=D_-/n\in(0,1]$. Let $c_*\in(0,1)$ be the root of $H(c)=\log2$, where
$H(c)=(1+c)\log(1+c)-c\log c$. Then $L_-/d\to c_*$, $c_*>1/4$, and
\begin{equation}
\liminf_{d\to\infty}y_d\ \ge\ \frac{c_*}{1+c_*}\ >\ \frac15,
\qquad
\frac{\tau_{>L_-}}{n\,\lambda_{d,L_-}}\ \longrightarrow\ 0 .
\label{eq:lastblock}
\end{equation}
\end{lemma}

\begin{proof}
$H$ is continuous and strictly increasing on $(0,\infty)$ ($H'(c)=\log\frac{1+c}c>0$), with $H(0^+)=0$ and
$H(1/4)=\frac54\log5-2\log2<\log2$ (numerically $0.6255<0.6931$), so $c_*$ exists and $c_*>1/4$.
By Stirling's formula, $\frac1d\log\binom{d+cd}{cd}\to H(c)$ uniformly for $c$ in compact subsets of
$(0,\infty)$. Fix $c<c_*$: by Lemma~\ref{lem:mult}, $\sum_{s\le cd}N_{d,s}\le(cd+1)\binom{d+cd}{cd}
=e^{d(H(c)+o(1))}<2^d$ for large $d$, so $L_-\ge cd$. Fix $c\in(c_*,1)$: $N_{d,cd}\ge\frac34\binom{d+cd-1}{cd}
=e^{d(H(c)+o(1))}>2^d$ for large $d$, so $L_-<cd$. Hence $L_-/d\to c_*$, and by Lemma~\ref{lem:mult}
$\rho_{L_-}=\frac{(d+2L_-)(d+L_--2)}{(L_-+1)(d+2L_--2)}\to\frac{1+c_*}{c_*}=:r_*$, with $r_*<5$ because
$c_*>1/4$.

For the first claim in \eqref{eq:lastblock}: $n<D_-+N_{L_-+1}$ by maximality of $L_-$, so
$y_d>\big(1+N_{L_-+1}/D_-\big)^{-1}$. Fix an integer $J$. For $1\le j\le J$ and $d$ large,
$\rho_{L_--j}=\rho_{L_-}\prod_{i=1}^{j}\frac{\rho_{L_--i}}{\rho_{L_--i+1}}\le\rho_{L_-}(1+1/(L_--J))^{2J}$,
because, from the product form of $\rho_s$ in Lemma~\ref{lem:mult},
$\frac{\rho_{s-1}}{\rho_s}=\frac{d+s-3}{d+s-2}\cdot\frac{s+1}{s}\cdot\frac{(d+2s-2)^2}{(d+2s)(d+2s-4)}
\le\big(1+\tfrac1s\big)\big(1+\tfrac{4}{(d+2s)(d+2s-4)}\big)\le\big(1+\tfrac1s\big)^2$ for $s\ge2$.
Writing $\bar\rho=\rho_{L_-}(1+1/(L_--J))^{2J}$,
$D_-\ge\sum_{j=0}^{J}N_{L_--j}\ge N_{L_-}\sum_{j=0}^{J}\bar\rho^{-j}$, while $N_{L_-+1}=\rho_{L_-}N_{L_-}$.
Since $L_-\to\infty$, $\bar\rho\to r_*$, and therefore
$\liminf_dy_d\ge\big(1+r_*\big/\textstyle\sum_{j=0}^J r_*^{-j}\big)^{-1}$ for every $J$; letting
$J\to\infty$, $\liminf_dy_d\ge\big(1+r_*(1-1/r_*)\big)^{-1}=1/r_*=c_*/(1+c_*)>1/5$.

For the second claim, by Lemma~\ref{lem:mult}, for $s\ge L_-\ge d/10$ (large $d$),
$\rho_s\le2\,\frac{d+s-2}{s+1}\le22$, and by Lemma~\ref{lem:mercer}, $\lambda_{s+1}/\lambda_s\le\gamma/d$.
Hence for $j\ge1$, with $L=L_-$,
$\frac{\tau_{L+j}}{n\lambda_L}=\frac{N_{L+j}}{n}\cdot\frac{\lambda_{L+j}}{\lambda_L}
\le\frac{N_L}{n}\,(22\gamma/d)^j\le(22\gamma/d)^j$, using $N_L\le D_-\le n$. Summing over $j\ge1$ gives
$\tau_{>L}/(n\lambda_L)\le 44\gamma/d$ for $d>44\gamma$.
\end{proof}

\begin{proof}[Proof of Theorem~\ref{thm:T3}]
Since $\gamma=h^{-2}$ is fixed, $d>\gamma$ for all large $d$, so Lemmas~\ref{lem:mercer}
and~\ref{lem:leverage} apply. Take $B=\{0,\dots,L_-\}$ in Lemma~\ref{lem:reduction}, so $D_B=D_-$ and
$T_B=T_{\le L_-}$. Taking
expectations in \eqref{eq:blockreduction} and applying Jensen's inequality, $\E[T^2]\ge(\E T)^2$, then
Lemma~\ref{lem:leverage},
\[
\frac{\E\,\mathcal E_{2^d,d}}{\sigma^2}=\E\frac{V(X)}{\sigma^2}\ \ge\ \frac{(\E T_{\le L_-})^2}{nD_-}
\ \ge\ \frac{nD_-}{\big(D_-+n-1+\tau_{>L_-}/\lambda_{d,L_-}\big)^2}
=\frac{y_d}{\big(1+y_d-\tfrac1n+\eta_d\big)^2},
\]
with $\eta_d=\tau_{>L_-}/(n\lambda_{d,L_-})\to0$ by Lemma~\ref{lem:window}. The function
$g(y)=y/(1+y)^2$ is increasing on $(0,1]$, since $g'(y)=(1-y)/(1+y)^3\ge0$. Given $\epsilon>0$,
Lemma~\ref{lem:window} gives $y_d\ge\frac15-\epsilon$ for all large $d$, and $y_d\le1$ always, so for large
$d$ the right side is at least $g(\tfrac15-\epsilon)-o(1)$. Letting $d\to\infty$ and then $\epsilon\to0$,
$\liminf_d\E\,\mathcal E_{2^d,d}/\sigma^2\ge g(1/5)=\frac{1/5}{(6/5)^2}=\frac5{36}$.
\end{proof}

\begin{remark}\label{rem:sharper}
The argument gives more than the stated constant. Lemma~\ref{lem:window} shows $\liminf y_d\ge1/r_*$
with $1/r_*=c_*/(1+c_*)=0.2271\ldots$, so the same proof yields the constant $g(1/r_*)=0.1509\ldots$;
the constant $5/36$ is the value at the round number $1/5$. The bound is an expectation statement from
start to finish: no concentration of the leverage is used, and none is needed. The certificate that the
proof evaluates, $(T_{\le L}/n)^2/(D_{\le L}/n)$, and the block dimension ratio $D_{\le L}/n$ are exactly the
quantities measured in Section~\ref{sec:t3}, where the one-window bound is compared with the realised
$V(X)/\sigma^2$ at $d\le13$. Running the same argument over several blocks simultaneously improves the
constant; the ten-window certificate $c_{10}=3882594969529/12500000000000=0.3106\ldots$ obtained in
\cite{PN24} is the value against which the measurements of Section~\ref{sec:t3} are also reported.
\end{remark}
